\documentclass[11pt, letterpaper]{article}

% --- Packages ---
\usepackage{amsmath}
\usepackage{graphicx}
\usepackage{amsfonts}
\usepackage{enumitem}
\usepackage{xcolor}
\usepackage{listings}
\usepackage{tcolorbox}
\usepackage[left=1in, right=1in, bottom=1in, top=3cm, headheight=2cm, headsep=0.5cm, footskip=1cm]{geometry}
\usepackage{fancyhdr}

% --- Code Listing Formatting ---
\definecolor{codegreen}{rgb}{0,0.6,0}
\definecolor{codegray}{rgb}{0.5,0.5,0.5}
\definecolor{codepurple}{rgb}{0.58,0,0.82}
\definecolor{backcolour}{rgb}{0.95,0.95,0.92}

\lstset{ 
    backgroundcolor=\color{backcolour},   
    commentstyle=\color{codegreen},
    keywordstyle=\color{blue},
    numberstyle=\tiny\color{codegray},
    stringstyle=\color{codepurple},
    basicstyle=\ttfamily\footnotesize,
    breakatwhitespace=false,         
    breaklines=true,                 
    captionpos=b,                    
    keepspaces=true,                 
    numbers=left,                    
    numbersep=5pt,                  
    showspaces=false,                
    showstringspaces=false,
    showtabs=false,                  
    tabsize=2,
    language=C++
}

% --- Header Configuration ---
\pagestyle{fancy} 
\fancyhf{} 
\renewcommand{\headrulewidth}{0.4pt} 
\lhead{\raisebox{-0.5cm}{\textbf{Universidad de Antioquia}}}
\rhead{\large \textbf{Tech Brief: GMAT Solver Sequences} \\ \normalsize Spaceflight Dynamics 2026-1} 
\cfoot{\thepage} 

\begin{document}

\begin{center}
    {\Large \textbf{Connecting Optimizers and Propagators in GMAT}} \\
    \vspace{0.2cm}
    {\large \textit{The Solver Control Sequence Architecture}} \\
    \vspace{0.4cm}
    \textit{Prof. Juan | Aerospace Engineering Program}
\end{center}
\hrule
\vspace{0.5cm}

\section*{The Architecture of a Solver Loop}

In GMAT, optimization is not a static calculation; it is a dynamic feedback loop. The \texttt{Optimize} command wraps around the \texttt{Propagate} command to form a \textbf{Solver Control Sequence}. 

The \texttt{Vary} command acts as the ``input'' to the physics engine, and the \texttt{NonLinearConstraint} and \texttt{Minimize} commands evaluate the ``output'' generated by the \texttt{Propagate} command.

\subsection*{Example: Low-Thrust Altitude Raising}

In this scenario, a spacecraft uses an electric thruster to raise its altitude. We want to find the exact \texttt{BurnDuration} that achieves a 12,000 km radius. Crucially, we want to minimize the fuel used.

Notice how the variable created by the Optimizer (\texttt{BurnDuration}) is fed directly into the Propagator's stopping condition.

\begin{lstlisting}[caption=GMAT Script: Connecting Optimize to Propagate]
% 1. Create a variable for the optimizer to manipulate
Create Variable BurnDuration;
Create Variable FuelCost;
BurnDuration = 1000; % Initial Guess in seconds

% 2. Open the Solver Control Sequence using the Yukon SQP Optimizer
Optimize Yukon1 {SolveMode = Solve, ExitMode = SaveAndContinue};

    % 3. VARY: Give the optimizer control over the BurnDuration
    Vary Yukon1(BurnDuration = BurnDuration, {Perturbation = 0.0001, MaxStep = 500.0, Lower = 0.0, Upper = 10000.0});
    
    % 4. ACTIVATE HARDWARE: Turn on the continuous thrust
    BeginFiniteBurn ElectricBurn(EP_SC);
    
    % 5. PROPAGATE (THE CONNECTION): The physics engine runs. 
    % It stops precisely when ElapsedSecs hits the Optimizer's current guess for BurnDuration.
    Propagate DefaultProp(EP_SC) {EP_SC.ElapsedSecs = BurnDuration};
    
    % 6. DEACTIVATE HARDWARE: Turn off the engine
    EndFiniteBurn ElectricBurn(EP_SC);
    
    % 7. CONSTRAINTS: Evaluate the state of the spacecraft exactly where the Propagate command stopped.
    % We force the optimizer to seek an exact radial distance of 12,000 km.
    NonLinearConstraint Yukon1(EP_SC.Earth.RMAG = 12000.0);
    
    % 8. MINIMIZE (OBJECTIVE): Define the cost function. 
    % We want to minimize the total fuel mass consumed during the propagation.
    FuelCost = EP_SC.FuelTank1.FuelMass;
    Minimize Yukon1(FuelCost);

% 9. END LOOP: Calculate gradients, update BurnDuration, and repeat Propagate.
EndOptimize;
\end{lstlisting}

\begin{tcolorbox}[colback=blue!5!white,colframe=blue!75!black,title=Key Takeaway for Students]
The \texttt{Propagate} command is the bridge between the mathematical guess and the physical reality. If you place a \texttt{NonLinearConstraint} before a \texttt{Propagate} command, the optimizer will evaluate the \textit{initial} state of the spacecraft, effectively doing nothing. The evaluation commands (\texttt{Constraint} and \texttt{Minimize}) must always be placed \textbf{after} the \texttt{Propagate} command so they can measure the physical consequences of the maneuver.
\end{tcolorbox}

\end{document}